\documentclass[12pt]{article}

\usepackage[doublespacing]{setspace}
\usepackage[
  letterpaper,
  left=1.5in,
  right=1in,
  top=1in,
  bottom=1in,
  footskip=0.5in
]{geometry}
\usepackage{amsmath,amssymb}
\usepackage{array,booktabs,tabularx}
\usepackage{enumitem}
\usepackage{float}
\usepackage[hidelinks]{hyperref}

\newcolumntype{L}[1]{>{\raggedright\arraybackslash}p{#1}}
\newcolumntype{Y}{>{\raggedright\arraybackslash}X}
\setlist[itemize]{leftmargin=0.28in,itemsep=0.35em,topsep=0.35em}
\setlength{\parindent}{0.5in}
\setlength{\parskip}{0pt}
\setlength{\emergencystretch}{2em}
\setcounter{secnumdepth}{0}

\makeatletter
\def\ps@plain{%
  \let\@mkboth\@gobbletwo
  \let\@oddhead\@empty
  \let\@evenhead\@empty
  \def\@oddfoot{\reset@font\hspace*{-0.5in}\hfil\thepage\hfil}%
  \let\@evenfoot\@oddfoot
}
\makeatother
\pagestyle{plain}

\hypersetup{
  pdftitle={Errata: A Low-Memory Spectral-Correlation Analyzer for Digital QAM-SRRC Waveforms},
  pdfauthor={Dylan Jacob Gormley},
  pdfsubject={Corrections to the May 2021 master's thesis}
}

\begin{document}

\hypersetup{pageanchor=false}
\begin{titlepage}
\thispagestyle{empty}
\begin{center}
\vspace*{0.9in}
{\large ERRATA\par}
\vspace{0.55in}
{\large A LOW-MEMORY SPECTRAL-CORRELATION ANALYZER\par}
\vspace{0.08in}
{\large FOR DIGITAL QAM-SRRC WAVEFORMS\par}
\vfill
{\large DYLAN JACOB GORMLEY\par}
\vfill
Original thesis submitted in partial fulfillment of the requirements for the degree of\par
MASTER OF SCIENCE IN ELECTRICAL ENGINEERING\par
at\par
CLEVELAND STATE UNIVERSITY\par
MAY 2021
\vfill
Errata revision: September 3, 2026
\end{center}
\end{titlepage}

\hypersetup{pageanchor=true}
\pagenumbering{arabic}

\begin{center}
ERRATA
\end{center}

\noindent\textbf{Use of this document.}
This errata records substantive corrections to the original May 2021 thesis. The corrected repository edition incorporates these corrections directly. Page references and Arabic table numbers below use the deposited thesis's printed pagination and labels (OhioLINK accession \texttt{csu1622636550863441}); the corrected edition formats Tables 10--13 as Tables X--XIII. Section, equation, figure, and table identifiers control if pagination differs in another copy. Minor typographical, grammatical, citation-style, and production changes are summarized at the end rather than itemized word by word.

\section{Abstract and scope}

\noindent\textbf{Affected location:} Abstract, p.~vii.

The abstract is clarified to state the scope of the characterization and implementation results precisely:

\begin{itemize}
  \item The thirteen characterization plots are single trials. Every plot has at least one threshold crossing, and eleven report the transmitted symbol-rate--center-frequency pair or pairs.
  \item Experiment VII reports a $1.00$-MBd rate for a $0.10$-MBd input. Experiment X reports a $0.04$-MHz center frequency for a $0.05$-MHz input.
  \item The work demonstrates an FPGA-targeted architecture through models, simulations, synthesis, and implementation resource counts. It does not report measurements from deployed or flight hardware.
  \item The implementation target is a ZedBoard XC7Z020\hspace{0pt}CLG484-1 using Vivado 2019.2.1.
\end{itemize}

\section{Random-process and spectral definitions}

\noindent\textbf{Affected locations:} Sections 2.1--2.2, pp.~4--14.

The following definitions replace the corresponding original descriptions.

\begin{itemize}
  \item A random process is a collection of random variables indexed by time. One outcome gives a sample path. For a real-valued process with a marginal density, its mean is the ensemble expectation at fixed time, for example
  \[
    \mu_X(t)=E\{X(t)\}=\int_{-\infty}^{\infty}x\,p_{X(t)}(x)\,dx
  \]
  when the marginal density exists. This ensemble average is not generally a time average along one record.

  \item A finite-second-moment process is wide-sense stationary when its mean is constant and its autocorrelation is invariant to a common time shift, so
  \[
    R_X(\tau)=E\!\left\{X\!\left(t+\frac{\tau}{2}\right)
      X^*\!\left(t-\frac{\tau}{2}\right)\right\}.
  \]
  A sample-path time average equals this ensemble quantity only under the relevant ergodicity condition.

  \item For a second-order cyclostationary process, the autocorrelation is periodic in absolute time. Its cyclic autocorrelation is the corresponding Fourier-series coefficient. A single-record time-average estimator requires the appropriate cycloergodicity assumption; it is not the population definition itself.

  \item Energy spectral density, power spectral density, and the periodogram are distinct. For an energy signal, the energy spectral density is $|X(f)|^2$. For a wide-sense stationary power process, the power spectral density is the Fourier transform of its autocorrelation. A periodogram is a finite-record estimate and is not a consistent power-spectral-density estimator without suitable averaging or smoothing. \mbox{Parseval's} relation equates total time-domain energy with integrated spectral energy; it does not identify a pointwise spectrum with power.

  \item The cyclic spectral density describes second-order periodic correlation. It does not specify a process's complete probability law. Its $\alpha=0$ slice is the ordinary power spectral density under the same conventions.

  \item The linear-modulation expression is
  \[
    X(t)=\sum_{k=-\infty}^{\infty}a_k p(t-kT_0-t_0)
    e^{j(2\pi f_c t+\theta_0)},
  \]
  where $T_0=1/R_s$. The phase offset belongs inside the exponential. Under the MATLAB convention used in the thesis, 2-QAM is an antipodal constellation equivalent to BPSK. The corrected QAM-SRRC derivation retains the symmetric-lag pulse product and the corresponding frequency shifts $P(f-f_c+\alpha/2)P^*(f-f_c-\alpha/2)$.
\end{itemize}

\section{Spectral-correlation analyzer and detector}

\noindent\textbf{Affected locations:} Sections 2.3--2.4, pp.~15--25; Sections 3.2--3.3.1, pp.~27--31; Chapter IV, pp.~49--55, including Tables 11--12 and Figures 35--37.

\begin{itemize}
  \item Gardner's baseline time-smoothing analyzer shifts the spectral components centered at $f+\alpha/2$ and $f-\alpha/2$ to baseband, low-pass filters them, correlates the two outputs, and takes the ordered time and bandwidth limits. It should not be described as an FFT of the cyclic autocorrelation.

  \item Koch's method first uses a thresholded Welch spectrum to form signal segments. Its cyclic-autocorrelation FFT applies an FFT to $|x[n]|^2$ to select candidate rates. For each selected rate, its implemented carrier detector sweeps a centered segment over frequency offsets, applies a rate-dependent low-pass filter, evaluates one cyclic-autocorrelation coefficient, and thresholds the normalized peak. A separate preliminary center-frequency preprocessor was not used in the reported system.

  \item For the serial architecture in this thesis, let
  \[
    z_k[n]=\sum_m h[m]y[n-m]
      e^{-j2\pi f_k(n-m)/F_s} .
  \]
  Its estimated zero-lag cyclic-autocorrelation coefficient and nonnegative test statistic are
  \[
    \widehat C_y(\alpha_v,f_k)=\frac{1}{N}\sum_{n=0}^{N-1}|z_k[n]|^2
      e^{-j2\pi\alpha_v n/F_s},
    \qquad
    T_y(v,k)=|\widehat C_y(\alpha_v,f_k)|^2.
  \]
  This is a channelized, filter-localized cyclic statistic, not literally a point value of the cyclic spectral density.

  \item The symbols $v$ and $k$ are array indices. The physical candidates are $\alpha_v$ and $f_k$ in hertz, and $F_s$ must appear in the discrete-time exponential. The Chapter IV plot coordinates are $\alpha_v/F_s$ and $f_k/F_s$ in cycles per sample.

  \item For the Chapter IV characterization, the plotted post-processing decision threshold is
  \[
    \gamma=\beta\widehat P_y^{\,2},\qquad \beta=10^{-4},
  \]
  where $\widehat P_y^{\,2}$ is the normalized sample-power-squared quantity that was mislabeled $\sigma_\eta^4$ in Table 12. The CFE testbench writes the statistic; the threshold overlay and decision are performed in post-processing. The value of $\beta$ is empirical rather than a calibrated probability of detection or false-alarm rate.

  \item A cyclostationary feature detector can reject ideal stationary noise at nonzero cycle frequency in the population model, but this does not make finite-record detection automatically immune to noise, uncertainty, leakage, or model mismatch. Adjacent grid crossings near one feature are correlated responses and should not automatically be counted as independent false alarms.

  \item The regenerated MATLAB illustrations use a constant threshold and independently impose the requested $E_b/N_0$ on both signals in the two-signal example. These explanatory figures do not alter the Chapter IV VHDL result plots.
\end{itemize}

\section{Architecture, memory, and numerical behavior}

\noindent\textbf{Affected locations:} Sections 3.1--3.3, pp.~27--48.

\begin{itemize}
  \item The cited strip spectral correlation analyzer intermediate array requires
  \[
    (2^{16})(32)(32\ \text{bits})=67{,}108{,}864\ \text{bits}
  \]
  for the stated $N=2^{16}$, $N_p=32$, and 32-bit complex-word assumptions. This exceeds the approximately $4.9$ Mbit of block RAM listed for the XC7Z020 in Xilinx DS190. The comparison is an order-of-magnitude architectural motivation, not a synthesis result for the streaming design.

  \item The low-memory architecture replaces a full-capture FFT array with bounded state, delay elements, lookup tables, and pipeline registers. At one sample per clock, a fully serial sweep over $N_v$ rate candidates and $N_k$ center-frequency candidates requires approximately $NN_vN_k$ sample clocks plus control and pipeline overhead.

  \item A $W$-bit numerically controlled oscillator represents a requested frequency with the nearest frequency-control word, giving quantum $F_s/2^W$. With $W=16$ and $F_s=1$ MHz, the hardware frequency resolution is $15.2588$ Hz. Experiment XI's 201 one-hertz-spaced requests therefore quantize to 14 distinct settings. The two discrete-time endpoints $-F_s/2$ and $+F_s/2$ map to the same oscillator setting.

  \item The thesis configuration uses \texttt{TAP\_COUNT=2} with the two-tap moving-average response $h=[0.5,0.5]$.

  \item The documented datapath uses a continuous-rate, fixed-latency streaming contract. During each record, \texttt{SampleInValid} is asserted on every clock and downstream \texttt{SampleOutReady} is held high. Sample-preserving stages advance data and their valid and final-sample tags together. Integrate-and-dump stages incorporate every accepted input into record state and emit one aligned result and final-sample tag per record. A fully elastic variant advances all sample-indexed state only on input acceptance and holds pending output data and tags until downstream acceptance.
\end{itemize}

\section{RTL numerical and interface requirements}

\noindent\textbf{Affected locations:} Sections 3.3.1--3.3.6 and 3.3.8--3.3.9, pp.~30--42 and 44--48.

The corrected edition states the following behavior as the RTL contract for the implementation.

\begin{itemize}
  \item For Q1.15 complex input components, the instantaneous power $P=I^2+Q^2$ is a nonnegative value in $[0,2]$ with a full-precision unsigned 32-bit representation and 30 fractional bits. The interface into the signed Fourier mixer preserves the magnitude bit by zero-extending into a wider signed word or by explicit rescaling and saturation. One conforming 16-bit mapping is saturated Q1.15 representation of $P/2$; this gives a one-half scale factor in the complex CAF coefficient and a one-quarter factor in its magnitude-squared statistic and matching threshold.

  \item The selected-frequency Fourier output is the normalized coefficient $\widehat X(f_0)=N^{-1}\sum_{n=0}^{N-1}X[n]e^{-j2\pi f_0n/F_s}$. Both component sums include sample $N-1$ before a single division by $N$, quantization, and emission.

  \item The two-tap FIR uses zero initial delay state for every candidate record. Within a record, the delay holds the preceding mixed sample. An inter-record clock with input valid deasserted inserts zero into the one-sample delay and inserts a deasserted valid tag for that corresponding delayed transaction; an earlier pipeline transaction may still emerge on the same clock. No state from one candidate enters the next.

  \item A signed $W$-bit input accumulated across $N$ samples uses $W+\lceil\log_2N\rceil$ bits for full-range operation. Thus, a 16-bit, 4096-sample configuration uses a 28-bit sum and a 12-bit normalization shift, while a 16-bit, 65,536-sample configuration uses a 32-bit sum and a 16-bit shift. For power-of-two record lengths, normalization uses a sign-aware rounded right shift. Other record lengths use a constant reciprocal or divider with guard bits before the same rounding and saturation stage.

  \item On each valid input, the accumulator first forms $S_{\mathrm{next}}=S+X[n]$. An ordinary sample stores that value. A sample carrying the final-sample marker at count $N-1$ emits the normalized $S_{\mathrm{next}}$, asserts aligned valid and final-sample tags, and clears the stored sum for the next record. An invalid input clock leaves the sample-indexed state unchanged and inserts deasserted valid and final-sample tags into the fixed-latency tag path; a previously accepted result may still emerge on that clock. This convention includes every record sample exactly once.
\end{itemize}

\section{Chapter IV experiments and results}

\noindent\textbf{Affected locations:} Chapter IV opening, p.~49; Tables 10--13, pp.~50--56; Figures 35--37, pp.~52--54.

Table 10 must state that $F_s=1$ MHz except in Experiment VII, where $F_s=10$ MHz, and that the single $E_b/N_0$ value in Experiments XII and XIII applies to both signals. In Table 11, the columns labeled $v$ and $k$ are requested physical grids and must instead be labeled $\mathcal R$ (MBd) and $\mathcal F$ (MHz). Figures 35--37 show the streaming-SCA test statistic, not a direct cyclic-spectral-density estimate.

\begin{center}
\singlespacing
\small
\renewcommand{\arraystretch}{1.18}
\begin{tabularx}{\textwidth}{@{}L{1.45in}L{2.20in}Y@{}}
\toprule
\textbf{Experiments} & $\boldsymbol{\mathcal R}$ \textbf{(MBd)} & $\boldsymbol{\mathcal F}$ \textbf{(MHz)} \\
\midrule
I--III, V, VI, VIII, IX & $[0.10:0.10:0.50]$ & $[-0.50:0.05:0.50]$ \\
IV & $[0.01:0.01:0.50]$ & $[-0.50:0.05:0.50]$ \\
VII & $[1.00:1.00:5.00]$ & $[-5.00:0.05:5.00]$ \\
X & $\{0.10\}$ & $[-0.50:0.01:0.50]$ \\
XI & $[0.0999:10^{-6}:0.1001]$ & $\{0.05\}$ \\
XII--XIII & $[0.05:0.05:0.50]$ & $[-0.50:0.05:0.50]$ \\
\bottomrule
\end{tabularx}
\end{center}

\begin{table}[H]
\centering
\singlespacing
\caption{Corrected interpretation of the Chapter IV result entries. Pairs are $(R_s,f_c)$ in (MBd, MHz).}
\small
\setlength{\tabcolsep}{4pt}
\begin{tabularx}{\textwidth}{@{}L{0.55in}L{1.20in}L{1.25in}Y@{}}
\toprule
\textbf{Expt.} & \textbf{Input pair(s)} & \textbf{Reported pair(s)} & \textbf{Correction or qualification} \\
\midrule
IV & $(0.01,0.05)$ & $(0.01,0.05)$ & The original result table printed $0.10$ MBd. The corrected nominal sample-processing time is $68.8128$ s. \\
VII & $(0.10,0.05)$ & $(1.00,0.05)$ & The original result table repeated the input rate. The configured receiver grid begins at $1.00$ MBd, so the selected rate is its first candidate. \\
VIII & $(0.10,0.05)$ & $(0.10,0.05)$ & The $0.430241$-s reported delay is independently consistent with the $0.430080$-s sample-processing lower bound; the difference is approximately $161\,\mu$s. \\
X & $(0.10,0.05)$ & $(0.10,0.04)$ & The selected center frequency is one 10-kHz search step below the input value. \\
XIII & $(0.10,0.05)$; $(0.25,0.10)$ & $(0.10,0.05)$; $(0.25,0.10)$ & The original result table printed $0.25$ MHz for the second center frequency. \\
\bottomrule
\end{tabularx}
\end{table}

For readability, the corrected edition expresses the Table 12 delays in seconds by dividing the original nanosecond values by $10^9$. These are serial-sweep simulation times, not per-sample detector latencies. Experiment VIII's $0.430241$-s result agrees with the $0.430080$-s nominal sample-processing time plus finite pipeline and control overhead. A single trial per condition characterizes the selected operating point but does not estimate detection or false-alarm probabilities.

In Table 13, the second row labeled Experiment VII must be Experiment VIII. The corrected table reports FPGA resource utilization by experiment for the ZedBoard XC7Z020CLG484-1 implementation target.

\section{Implementation target and toolchain}

\noindent\textbf{Affected locations:} Chapter III, pp.~26--48; Section 4.2, p.~55; Table 13, p.~56; Chapter V, pp.~57--58.

The implementation targets a ZedBoard with part XC7Z020CLG484-1 and uses Vivado 2019.2.1 with the default synthesis and implementation strategies. The component and experiment tables use this target and toolchain convention. The corrected edition presents the fixed-point formats, record-boundary behavior, accumulator sizing and normalization, frequency-control-word mapping, and continuous-rate streaming interface as the implementation contract used throughout the design and characterization.

\section{Conclusions and claims of support}

\noindent\textbf{Affected location:} Chapter V, pp.~57--58.

The conclusion describes an FPGA-targeted architecture, MATLAB models, VHDL simulations, and synthesized resource utilization. The work does not report deployed flight-hardware measurements. The serial processing time grows with the capture length and the number of candidate pairs. Koch's preprocessing can reduce the number of candidates presented to the serial sweep and requires spectral segmentation, filtering, block storage, an FFT, and validation.

Any future decimation must be preceded by an anti-alias low-pass filter and must preserve the occupied bandwidth and required cyclic frequencies. The factor is not generally $F_s/R_s$. The Chapter IV fixed-point VHDL characterization tests only 4-QAM-SRRC and 16-QAM-SRRC inputs. Passing other waveform samples through the same datapath does not establish useful features, thresholds, or reliable support for FSK, OOK, MSK, DSSS, DPSK, Manchester coding, or other pulse shapes.

\section{Editorial, figure, and reproducibility amendments}

The corrected edition also makes the following non-substantive or production-oriented changes:

\begin{itemize}
  \item It corrects spelling, grammar, punctuation, capitalization, units, mathematical notation, cross-references, table labels, and bibliographic metadata throughout.
  \item It regularizes front-matter order, margins, page-number placement, chapter and contents formatting, and the presentation of figures and tables.
  \item It repairs the MATLAB plotting scripts, uses fixed random seeds where stochastic examples are shown, corrects formulas and search grids, and regenerates the figures with a shared font, size, line-weight, and color specification.
  \item It plots the 2-, 4-, 8-, and 16-QAM-SRRC cyclic spectra under the minimum-distance-two convention on one vertical scale, making the $1{:}2{:}6{:}10$ symbol-variance scaling explicit, and aligns the full-CSD threshold plot's viewing axes with the other three-dimensional spectra.
  \item It replaces the conceptual SCA and RTL architecture and timing raster images with native TikZ drawings that use the corrected notation, specified signal order, and a shared typographic style.
  \item It applies conventional DSP and FPGA symbols to the functional diagrams, uses \texttt{tikz-timing} for the timing diagrams, and regularizes the displayed interface aliases without changing the VHDL source identifiers.
  \item It adds a front-matter nomenclature of recurring symbols and abbreviations and removes the duplicate acronym appendix from the corrected edition.
  \item It replaces the stale Appendix A code transcription with the maintained MATLAB source included by the document build.
\end{itemize}

\section{Principal verification sources}

The corrections above were checked against the thesis source, experiment plots, component design records, executable models, and the following primary technical sources already represented in the corrected thesis bibliography:

\begin{itemize}
  \item W. A. Gardner, \textit{Introduction to Random Processes with Applications to Signals and Systems}, 2nd ed., McGraw-Hill, 1990.
  \item W. A. Gardner, A. Napolitano, and L. Paura, ``Cyclostationarity: Half a Century of Research,'' \textit{Signal Processing}, vol.~86, no.~4, pp.~639--697, 2006.\\
  \href{https://doi.org/10.1016/j.sigpro.2005.06.016}{doi:10.1016/j.sigpro.2005.06.016}.
  \item M. V. Koch and J. A. Downey, \textit{Interference Mitigation Using Cyclic Autocorrelation and Multi-Objective Optimization}, NASA/TM--2019-220226, 2019.
  \item Xilinx, Inc., \textit{Zynq-7000 SoC Data Sheet: Overview}, DS190, version 1.11.1, 2018.
\end{itemize}

\vfill
\begin{center}
\singlespacing
End of errata
\end{center}

\end{document}
